% =====================================================================
%  2 Related Work / Background
%  Sources: reports/Final Thesis Cross-Cultural Disagreement ... .md
%  (Notion export, main list), reports/key_points_thesis_proposal.md,
%  reports/technical_papers_sum.md (PASTA, FairSteer, political steering,
%  TDD), reports/paper_summaries.md, reports/v3_step1_explained.md (good
%  plain-language explanation of attention, heads and steering).
%  End each subsection with one sentence on how this thesis differs.
% =====================================================================
\section{Related Work}
\label{sec:related}

\subsection{Annotator disagreement and perspectivism}
\label{sec:related:disagreement}
% - Disagreement as signal: Denton et al. 2021 (whose ground truth),
%   Uma et al. 2021 (learning from disagreement survey), Plank 2022
%   (human label variation).
% - Cultural and moral correlates: D3CODE (Davani et al. 2024),
%   Davani et al. 2024 FAccT (care-high raters stricter; Perspective API
%   aligns with WEIRD values), Lee et al. 2024 (cross-cultural hate speech,
%   sarcasm and local topics drive disagreement).
% - Modelling groups: jury learning (Gordon et al. 2022), GRASP
%   (Prabhakaran et al. 2024, group association metrics), Maurer et al.
%   2026 (text x annotator interaction), Muhumuza et al. 2026 (HateXplain).
% - Demographics explain little variance: Hu & Collier 2024.
% How we differ: we test whether the group signal is reliable at the item
% level before asking a model to reproduce it (section 7).
\fillin{Disagreement, perspectivism, D3CODE and moral correlates.}
\citep{denton2021whose,uma2021learning,plank2022problem,davani2024d3code,davani2024disentangling,lee2024exploring,gordon2022jury,prabhakaran2024grasp}

\subsection{LLMs as offensiveness judges and persona prompting}
\label{sec:related:persona}
% - Overconfidence on contested items: Lu et al. 2025.
% - Persona effect is small and bounded by how much demographics explain:
%   Hu & Collier 2024. More attributes can hurt: Kamruzzaman et al. 2025.
% - Persona prompting barely changes LLM bias patterns: Giorgi 2025;
%   persona effects depend on the perspectivism setting: Sarumi et al. 2025.
% - LLMs fail to predict group differences in safety ratings: Movva et
%   al. 2024. Country prompts do not change GPT-4 labels: Lee et al. 2024.
% - Persona prompting can introduce bias: Salewski et al. 2023, Gupta et
%   al. 2024; anti-Muslim bias: Abid et al. 2021. Robust persona-aware
%   toxicity detection with ensembling: Atil et al. 2025.
% How we differ: content-free, length-matched controls; place-name
% controls; mechanism, not only behaviour.
\fillin{LLMs as judges, persona prompting, what is known about its effect.}
\citep{lu2025overconfident,hu2024quantifying,kamruzzaman2025demographic,giorgi2025human,sarumi2025impact,movva2024annotation,salewski2023incontext,gupta2024bias}

\subsection{Prompt sensitivity}
\label{sec:related:prompt}
% - Recency and majority-label bias, calibration: Zhao et al. 2021.
% - Position effects in long contexts: Liu et al. 2024.
% - Relevance for us: any inserted sentence moves the answer; suffix
%   placement reverses effects in Qwen; phrase length matters.
\fillin{Prompt-format effects and why a content-free control is needed.}
\citep{zhao2021calibrate,liu2024lost}

\subsection{Attention analysis and interpretability}
\label{sec:related:attention}
% - Attention as explanation debate: Jain & Wallace 2019; Wiegreffe &
%   Pinter 2019. Our profiling result (correlational candidates overlap
%   causal heads barely above chance) fits this debate.
% - Value-weighted attention (norm-based): Kobayashi et al. 2020.
% - Attention sinks: Xiao et al. 2024 (about 80% of attention goes to the
%   first token in our pilots).
% - Personas inside the model: Ghandeharioun et al. 2024 (layer-wise
%   persona representations). Prompt influence / token attribution:
%   Feng et al. 2024 (TDD).
\fillin{Attention analysis, its limits, sinks, value-weighted attention, persona mechanisms.}
\citep{jain2019attention,wiegreffe2019attention,kobayashi2020attention,xiao2024efficient,ghandeharioun2024whos,feng2024unveiling}

\subsection{Steering language models}
\label{sec:related:steering}
% - Attention steering: PASTA (Zhang et al. 2024): upweight attention to a
%   user-marked span in selected heads, heads chosen by profiling.
% - Activation steering: ActAdd (Turner et al. 2023), contrastive
%   activation addition (Rimsky et al. 2024).
% - Steering for bias and fairness: FairSteer (Li et al. 2025), political
%   views steering papers (see reports/technical_papers_sum.md, add bib
%   entries once you decide to cite them).
% How we differ: we steer a demographic span and judge success against
% real group ratings AND a constant-shift baseline.
\fillin{Attention steering (PASTA), activation steering, steering for bias.}
\citep{zhang2024pasta,turner2023activation,rimsky2024steering,li2025fairsteer}

\subsection{Positioning}
% One paragraph: the combination nobody has done (behavioural controls +
% mechanism + steering + reliability of the group signal) on one dataset.
\fillin{What this thesis adds relative to the work above.}
